% ED26-BM-RAT2021-10-6: complete independently worded source reconstruction.
\chapter{Adaptationism: A Meta-Normative Theory of Rationality}
\label{sec:p11-rationality-10-6}

Cosmides and Tooby ask which standards appropriately assess an evolved
information-processing mechanism. Their answer makes a theory of its
adaptive function prior to the choice of a performance benchmark.
An ancestral problem, the information then available and the consequences
of alternative responses constrain what would count as good design.
The chapter's five cases concern data format, need-sensitive risk,
activation cues, content-specialized social reasoning and cooperation
under uncertainty about repeated interaction
\cite[chap.~10.6, pp.~633--642]{CosmidesTooby2021Adaptationism}.

This is a meta-normative proposal: it concerns the justification and
applicability of standards, rather than furnishing one content-free rule
for every inference or choice. Its meaning of rationality is functional
good design relative to an adaptive problem. Whether that standard also
has epistemic, practical or moral authority is a further question.
The evolutionary explanation and that normative question must remain
distinct, even when one response satisfies both standards.

\section{The argument for an adaptive-function benchmark}
\label{sec:p11-rat106-argument}

The source begins with a contrast. Laboratory reasoning has often been
compared with logic, Bayesian probability, statistical decision theory,
falsification or expected utility, with departures classified as errors.
The authors associate the resulting picture with a heterogeneous
automatic System~1 and an effortful, deliberative System~2. They contrast
that picture with organisms' successes at language, perception and
social inference. Their comparison with the AI available when they wrote
is historical motivation, not a dated performance dataset or a claim
about every later AI system. Neither natural success nor a laboratory
error label by itself establishes the appropriate norm
\cite[pp.~633--634]{CosmidesTooby2021Adaptationism}.

The first argumentative step is goal relativity. The source contrasts
charcoal used to relieve hunger with charcoal used after poisoning,
direct travel with zigzagging to evade a predator, and talking alone
with rehearsing a role. The same coarse behavior can serve different
functions; one cannot evaluate its design without fixing the relevant
problem. These are conceptual illustrations, not medical instructions,
participant observations or proofs that every action has only one goal.
They challenge an unspecified benchmark, rather than abolish truth
preservation or probability laws within their proper domains.

The second step concerns the object of assessment. Modern scientific
inquiry, markets of anonymous agents and large-scale cooperation have
goals and information resources unlike those that selected ancestral
cognitive machinery. The source associates these historical changes
from the fifteenth century onward with increasing literacy, durable
data storage and cheaper communication and travel, enabling wider
comparison of hypotheses and cumulative knowledge, anonymous markets
and contact among communities with different moral conventions. This
is its historical framing, not a measured causal estimate of those
developments. Cosmides and Tooby argue that a standard developed
for such a modern task need not be a measure of the mechanism's evolved
design. This does not entail that the standard is useless for the modern
task, or that people ought never to learn it. The distinction is between
explaining the design that arose and recommending a procedure for a
present purpose \cite[p.~634]{CosmidesTooby2021Adaptationism}.

\begin{hbdefinition}{def:p11-rat106-function}{Adaptive problem, adaptation and selected function}
In this source, an adaptive problem is a recurrent ancestral condition or
causal relationship that creates a reproductive opportunity or obstacle,
and that a possible property could exploit to improve reproduction of
its bearers or relatives carrying the relevant inherited design.
Recurrence alone is insufficient: the reproductive relationship matters.

An adaptation is an organized design whose features were selected because
they solved such a problem. Its selected function is that problem-solving
effect, not every effect that is presently beneficial. A computational
adaptation processes information to produce solutions. Heart pumping,
liver detoxification and retinal transduction illustrate the distinction
between a structure, its causal organization and its selected function.
The definition does not establish that any particular psychological
trait is an adaptation; that classification requires evidence.
\end{hbdefinition}

The source's selection narrative supplies the explanatory premise.
Heritable changes that impair machinery tend to reduce its bearers'
reproductive contribution; changes that improve it can leave more copies
of the design. Iteration can accumulate coordinated functional features.
The authors describe negative and positive feedback and the emergence
of species-typical organization. This is a qualitative account, not a
population-genetic theorem that every beneficial mutation fixes, every
common trait is adaptive or every response achieves an optimum.
Drift, alternatives absent from a population, changing environments,
tradeoffs and historically constrained designs prevent those stronger
inferences. Selection on reproduction also does not imply a person's
conscious goal is reproductive success
\cite[pp.~634--635]{CosmidesTooby2021Adaptationism}.

\begin{hbmethod}{def:p11-rat106-task-analysis}{The source's adaptive task analysis}
Identify a candidate recurrent adaptive problem before describing the
mechanism as well designed. Reconstruct the ancestral information,
uncertainty, alternatives and relevant consequences of success and error.
Ask which representations, inference operations and decision procedures
could exploit that information to solve the problem. Derive discriminating
predictions about inputs, outputs, activation and failure; then test them
against observations and competing mechanisms.

The proposed normative bridge is conditional: if rationality of a
computational adaptation means good design for its selected function,
and the hypothesized adaptive problem really explains the mechanism's
design, task analysis supplies a nonarbitrary functional benchmark.
Its substantive work lies in identifying the function and deriving
what improves its solution. Merely declaring an observed response
adaptive supplies neither step.
\end{hbmethod}

This reconstructs the argument at printed 635. The goal-relative premise
motivates asking for a function; the evolutionary premise explains why
selected functions are relevant to evolved design; task analysis connects
the function to testable standards. The conclusion concerns design in
that respect. It does not follow solely from a behavior's prevalence,
nor yield an unrestricted ought. An objection can therefore target the
selected-function hypothesis, the derived performance model, or the
claim that functional goodness exhausts rationality. Changing a name
does not answer any of those objections.

\section{Case one: probability information in a usable format}
\label{sec:p11-rat106-frequencies}

The source compares color constancy in ordinary daylight with failure
under the unusual spectrum of sodium-vapor light. The purpose is to
illustrate input mismatch: failure in a novel input environment need
not reveal poor design for the environment that selected the mechanism.
It then makes the analogous claim about single-event probabilities and
normalized percentages versus naturally sampled event counts.
Probability theory is not rejected. The proposed difference is which
representation makes a valid computation accessible
\cite[p.~636]{CosmidesTooby2021Adaptationism}.

The source also contrasts the favorable Bayesian descriptions of animal
foraging in Real 1991 and Stephens and Krebs 1986 with unfavorable
human laboratory descriptions. That comparison motivates investigation
of ecological inputs; it supplies neither one common cross-species task
nor a quantitative ranking of animal and human statistical competence.

\begin{hbexample}{def:p11-rat106-orchard}{The source's orchard count}
An orchard has 125 trees. Five are apple trees, giving an apple frequency
of 4 percent. Among the 50 trees with red fruit encountered in the
example, five are apple and 45 are cherry. Relative to this specified
finite count, the share of red-fruit trees that are apple is
\[
                         \frac{5}{5+45}=\frac{1}{10}.
\]
The denominator is the red-fruit set, not all 125 trees. Thus the low
apple base rate is already embodied in the five apple cases within
that count. It has not been ignored as evidence.

This is a worked conditional frequency, not a newly observed orchard
or a study of people's performance. Applying it to an unknown population
requires an appropriate sampling model. It cannot license ignoring a
separately relevant base rate when supplied conditional percentages
were normalized within different groups.
\end{hbexample}

Cosmides and Tooby attribute improved Bayesian-task performance with
natural frequencies to Gigerenzer and Hoffrage 1995/1999 and their own
1996 work; Kleiter's natural-sampling analysis supplies the representation
connection. They also describe such formats as ancestrally available
and point to modern risk communication. These are the selected chapter's
survey claims, not an independent estimate that all people automatically
encode every frequency or succeed on all Bayesian tasks. The chapter
does not provide the studies' recruitment, sample denominators,
assignment procedures, task batteries, effect intervals or transfer
results. Those details are required before estimating a format effect
for a named population or prescribing a particular clinical aid.
The formal count and the adaptation hypothesis have different evidential
standing \cite[p.~636]{CosmidesTooby2021Adaptationism}.

\section{Case two: risky choice depends on the relevant need}
\label{sec:p11-rat106-risk}

The disease example offers a sure saving of 200 of 600 lives, or a
lottery saving all 600 with probability one third and none otherwise.
Describing the same outcomes as 400 certain deaths, or zero deaths
with probability one third and 600 otherwise, changes the frame but
not the stipulated distributions. The chapter reports approximately
70 percent choosing the sure gain-framed program and approximately
80 percent choosing the risky loss-framed program, attributed to
Tversky and Kahneman 1981. These are chapter-level rounded reports;
the chapter supplies no sample denominators, recruitment frame,
assignment details or intervals. They are not universal rates, an
individual longitudinal preference cycle, or a replication synthesis
\cite[p.~636]{CosmidesTooby2021Adaptationism}.

\begin{hbdefinition}{def:p11-rat106-disease-model}{The lottery and a threshold criterion}
For this formal example, let \(A\) be the certain outcome of 200 lives
saved and \(B\) the lottery with 600 saved with probability \(1/3\)
and zero with probability \(2/3\). Let \(u\) be a real-valued utility
on these outcomes. Expected-utility comparison uses \(u(200)\)
and \(u(600)/3+2u(0)/3\).

For a different stated goal, let \(k\) be the minimum number of lives
that must be saved, with \(0\leq k\leq600\). The threshold criterion
maximizes the probability of saving at least \(k\). It uses the
indicator utility \(u_k(x)=1\) when \(x\geq k\), and zero otherwise.
This choice of criterion is an assumption, not a result established
by the number of lives or by evolutionary history.
\end{hbdefinition}

\begin{hbtheorem}{thm:p11-rat106-lottery}{proved=annex:p11.rat106-lottery}{Equal means, different utilities and the need threshold}
Both \(A\) and \(B\) have expected lives saved 200. They are
expected-utility indifferent exactly when
\(u(200)=u(600)/3+2u(0)/3\); affine utility implies this equality,
but equal means alone do not. Under the threshold criterion both
options tie at \(k=0\); \(A\) is strictly better when
\(0<k\leq200\); \(B\) is strictly better when
\(200<k\leq600\). Changing between saved and died descriptions
does not change either comparison when the consequence utility and
the need threshold are held fixed.
\end{hbtheorem}

The source's unqualified claim that expected utility predicts
indifference because the expected values coincide needs the utility
qualification just supplied. A nonlinear utility can prefer either
lottery without violating expected-utility representation. Similarly,
a choice reversal between two frames is a failure of description
invariance under a common interpretation, not by itself a proof of
intransitivity of one fixed preference relation. A demonstrated
transitivity violation needs its own comparisons. These corrections
preserve the source's important question about benchmark adequacy
without attributing a false implication to decision theory.

The source's risk-sensitive foraging illustration has two patches with
mean daily yield about 200 seeds but different variability. A bird
needing 100 seeds is contrasted with one needing 300 to survive until
the next day. Minimizing the chance of falling short can favor a
low-variance option in the first situation and a high-variance option
in the second. Exact choices require the full distributions: means,
variances and need alone do not determine every lower-tail probability.
For instance a patch paying exactly 200 cannot meet a need of 300,
whereas a varying patch with positive probability above 300 can.
This makes the proposed criterion operational without pretending the
chapter's sketch supplies every patch distribution
\cite[p.~637]{CosmidesTooby2021Adaptationism}.

The authors attribute human need manipulations and framing results to
Mishra and Fiddick 2012 and Rode and colleagues 1999. They propose that
frames change the implicitly adopted need and explain why some choices
called loss aversion depend on context. Rode's reported preference for
an ambiguous option when it is the lower-variance alternative is used
against unconditional ambiguity avoidance. These are bounded survey
reports of those task findings, not demonstrations that loss aversion
or ambiguity aversion never exists. No participants, measures beyond
the described choices, uncertainty intervals or primary-study designs
are supplied in this chapter. A need-sensitive explanation competes
with preference, interpretation and framing accounts; performance
consistent with it does not by itself identify an ancestral mechanism.

\section{Case three: activating a mechanism in its domain}
\label{sec:p11-rat106-activation}

\begin{hbdefinition}{def:p11-rat106-domains}{Proper and actual domains}
Following the distinction attributed to Sperber, the proper domain is
the information a mechanism was selected to process. Its actual domain
consists of the inputs that do activate it, including novel inputs
resembling the selected cues. The source's ground-squirrel illustration
contrasts an approaching hawk with a drone shaped like one: the latter
can elicit an alarm without being the ancestral danger. The example
illustrates the distinction; it is not a newly documented animal study.
\end{hbdefinition}

Failure to produce an inference in an uncontextualized task leaves two
possibilities open: a relevant procedure may be absent, or its activation
conditions may not be met. The source's plumber example contrasts a
bare conditional about repair and a dry floor with an argument in which
the plumber's repair claim is challenged. Wetness the next morning
supports denial of repair under the stipulated conditional.

\begin{hbtheorem}{thm:p11-rat106-conditional}{proved=annex:p11.rat106-conditional}{Material conditional violations and the plumber inference}
For classical propositional truth values, \(P\to Q\) is false
exactly when \(P\) is true and \(Q\) is false. Consequently
\(P\to Q\) together with \(\neg Q\) entails \(\neg P\).
In a four-case information-search task with exactly one known side
per case and one case of each known type \(P,\neg P,Q,\neg Q\),
the cases known to satisfy \(P\) and those known to
satisfy \(\neg Q\) are the only types that could conceal a
violation; known \(\neg P\) or known \(Q\) cannot conceal one.
This concerns possible violations of a material conditional, not
every pragmatic or deontic interpretation of an English conditional.
\end{hbtheorem}

The chapter surveys approximately 75 percent failure on context-free
modus-tollens tasks, citing Rips 1994 and Wason and Johnson-Laird 1972.
It supplies no particular sample, condition allocation, item denominator
or interval for that percentage. The plumber contrast is an illustration
of the activation hypothesis, not a controlled estimate that argument
raises every person's success rate
\cite[pp.~637--638]{CosmidesTooby2021Adaptationism}.

The proposed evolutionary rationale comes from the communication
account attributed to Krebs and Dawkins and the argumentative account
of Mercier and Sperber. Stable signaling can benefit senders and
receivers while their interests diverge on particular messages.
Deceptive signals can favor discrimination by receivers, followed by
harder-to-detect deception by senders. Honest signaling, where interests
converge, creates a different design problem, including resistance to
cheap imitation. The source cites Higham for costly honest signaling.
This narrative assumes the relevant selection/payoff relationships;
it is not a theorem that every signal benefits its current receiver or
that all communication is deception.

Language does not automatically make lying more expensive in words
than truth telling. In the argumentative proposal, epistemic vigilance
therefore evaluates a message against other statements and known facts.
The source describes a meta-representational format that keeps a claim
under assessment apart from accepted semantic knowledge. Disagreement
can reflect misinformation or manipulation, either of which motivates
checking; a persuader must also construct an effective argument.
This supplies a functional hypothesis for logical reasoning's activation,
not proof that argument establishes truth or that persuasion is justified.

The proposal does not predict flawless logic in every argument.
Reasoning for a desired conclusion can produce confirmation bias,
whereas a receiver's vigilance can demand scrutiny. The source takes
that contrast to fit an argumentative function. But functional
effectiveness at persuading someone differs from evidential support
for the conclusion. The psychological mechanism, evolved function,
and epistemic evaluation are three distinct claims
\cite[p.~638]{CosmidesTooby2021Adaptationism}.

\section{Case four: cheater detection is content-sensitive}
\label{sec:p11-rat106-exchange}

Domain-general computation exists in the source's account, but is not
its entire proposed architecture. Belief and desire help predict a
person rather than a rock; likewise exchange reasoning can use
representations of agent-specific benefits and requirements. The
adaptive problem is conditional cooperation: obtaining gains from
cooperators without repeatedly paying uncompensated costs to cheaters.
The relevant cheater intentionally accepts a benefit without meeting
the supplier's requirement, not merely someone whose innocent error
produces the same overt noncompliance
\cite[pp.~638--640]{CosmidesTooby2021Adaptationism}.

The authors appeal to evolutionary game models associated with Trivers,
Axelrod and Maynard Smith. Strategies can be contingent rules without
conscious backward induction; payoffs must be related to reproduction
and inherited strategies before they represent fitness. In models where
unconditional providers face exploiters and uncompensated giving lowers
reproductive payoffs, indiscriminate cooperation can lose to cheating.
That conditional statement is not a general impossibility of generosity:
relatedness, reputation, structure, different payoff relations or
transmission mechanisms change the problem. Successful conditional
cooperation requires relevant information about exploitation, but does
not itself prove a particular human neural implementation.

The chapter says that the exchange-specific hypotheses were constructed
before the experiments' data were collected. This is its account of
the research sequence, not an independently reconstructed preregistration
history. The distinction matters because fitting an adaptive explanation
to a known response pattern is a weaker test than predicting a new
pattern that competing explanations would not predict.

The source's Wason illustration asks which partly described birds
could violate the rule that a swan is white: investigate the swan's
color and the black bird's species, not the already white bird.
Theorem~\ref{thm:p11-rat106-conditional} supplies the formal reason.
The exchange version concerns borrowing a car conditional on refilling
its tank: inspect the borrower and the person who did not refill.
The chapter reports fewer than 25 percent choosing exactly the
logically relevant cases on indicative tasks, contrasted with 65--80
percent finding violations on exchange tasks. These are its aggregate
survey percentages, not identified denominators or a matched treatment
effect; their samples, item versions, recruitment and uncertainty are
not reported in the selected chapter
\cite[pp.~639--640]{CosmidesTooby2021Adaptationism}.

\begin{hbexample}{def:p11-rat106-exchange-direction}{The content goal and the displayed conditional can diverge}
In the source's car agreement, someone receiving the benefit without
meeting the requirement is also a material-conditional counterexample.
Agreement of answers there cannot distinguish general falsification
from benefit/requirement tracking.

For an explicit editorial contrast, suppose the same agreement is
being policed but the displayed descriptive rule is instead: everyone
who filled the tank borrowed the car. Its material-conditional
counterexample is a refiller who did not borrow. The agreement's cheater
remains a borrower who did not refill. These are different individuals
and different search questions. The original agreement does not assert
the converse descriptive rule. A task that elicits cheater search under
such a manipulation must establish what participants understood.
\end{hbexample}

Cosmides and Tooby describe a specialized exchange system with
representations of benefits, agent-relative requirements, obligation,
entitlement and cheating. Their claim is that it tracks cheaters even
when the resulting selection differs from content-free logical
violation search, and that innocent errors do not activate the same
search. They also report functional and neural dissociation from
precautionary-rule reasoning and refer to their 2015 review and 2008
deontic comparison. These are explicitly the chapter's survey boundaries;
it contains no lesion/imaging sample, scanner contrast, effect estimate,
design details or independent replication record. Improved performance,
selective impairment and neural contrasts have different inferential
requirements. None follows merely from writing a content-rich vocabulary.
The proposed design is evaluated relative to exchange's adaptive task,
not licensed as logically valid or morally right in every context.

\section{Case five: uncertainty about whether an encounter repeats}
\label{sec:p11-rat106-reciprocity}

The chapter contrasts economic models with observed generosity,
cooperation, trust and punishment in anonymous laboratory dictator,
prisoner's-dilemma, bargaining, trust and public-goods tasks. Its
immediate-payoff prediction presupposes the relevant payoff preferences
and game, not simply rationality alone. Other-regarding utility, repeated
interaction, reputation, relatedness and beliefs about the instructions
must not be excluded merely by calling a participant rational.
The cited Hoffman and colleagues and Smith accounts are source-survey
comparisons; this chapter gives no common human sample or cross-game
effect estimate \cite[p.~640]{CosmidesTooby2021Adaptationism}.

The actual information-processing question is whether the encounter
will repeat. The authors invoke ancestral small-group life and a
Krasnow and colleagues 2013 model in which meeting now informs the
chance of meeting later. They distinguish falsely treating a one-shot
encounter as repeated from missing a genuinely repeated relationship.
Paying one unnecessary cooperative cost can be less costly than
forgoing a long series of gains. This is a proposed asymmetric-error
account, not evidence that every one-shot cost is smaller or that
ancestral base rates are known numerically.

\subsection{The actual dependent simulation}

Delton and colleagues' primary study supplies the modeled task and
design, not a participant experiment. Each generation has 500 simulated
agents, randomly paired independently of genotype and randomly assigned
to one-shot or indefinitely repeated dyads. Each agent observes its
partner's noisy cue. Conditional on the dyad type, cues are independent
normal draws with standard deviation one, means \(-d/2\) for one-shot
and \(d/2\) for repeated interactions. The tested distances are
\(d=1,2,3\). No third-party reputation effects are included; partners
can identify one another in repeated interaction
\cite[pp.~13336--13337; SI~pp.~1--2]{DeltonEtAl2011Reciprocity}.

On a round, cooperation delivers benefit \(b\) at personal cost \(c\),
with \(b>c>0\). The simulations fix \(c=1\), baseline fitness 10
and \(b\in\{1.5,2,3,4,5,6,7,8,9,10\}\). One-shot shares are
10, 30, 50, 70 or 90 percent. In repeated dyads each next round occurs
independently with probability \(w\in\{0.5,0.8,0.9,0.95,0.99\}\),
giving mean lengths 2, 5, 10, 20 or 100. These are experimental model
parameters, not measured ancestral frequencies. Payoff-proportional
asexual reproduction fills the next generation; offspring inherit
regulatory values subject to mutation. Runs last 10,000 generations.
Mutation probability is 5 percent, with zero-mean normal increments
of standard deviation 0.025; probability-valued traits remain in
\([0,1]\). The first generation receives mutation to generate variation.

In SS1, the belief component knows the cue distributions and base rates
and uses Bayesian updating. A dyad is classified by which posterior
type probability is larger; a classification as one-shot is not certainty.
Two inherited traits give cooperation probabilities conditional on
each belief classification. Initially these are near zero for a
one-shot belief and near one for a repeated belief, with starting
mutation supplying variation. Cooperation implements TIT-for-TAT,
whereas defection implements always-defect. SS2 instead fixes the
mapping from classification to those actions and lets a cue threshold
evolve. Larger cues indicate repetition; lower thresholds therefore
make cooperation easier to activate. These are two restricted
architectures, not an exhaustive competition among all strategies or
a determination of human cognition's actual architecture.

Each architecture has 750 parameter combinations and four runs per
combination, hence 3,000 runs. Repeating both with a 5-percent chance
of accidental defection when intending cooperation, using GRIM's
permanent withdrawal after an observed defection, gives 12,000 runs
in all. Recorded summaries average the final 500 generations; supplementary
figures also average over cue distances. Random draws and reproduction
give simulation variability, not uncertainty about a sampled human
population \cite[SI~pp.~1--3]{DeltonEtAl2011Reciprocity}.

For a reported SS1 example with one-shot share one half, \(d=2\)
and \(b/c=3\), evolved cooperation given a one-shot classification
is approximately 87 percent for mean repeated length five and 96
percent for length ten. The primary reports approximately 16 percent
of those classifications are actually repeated in this setting.
With \(d=3\) and mean length ten, the corresponding examples are
47 percent cooperation and 7 percent classification error. These
are different settings, not one common effect or a universal 70--90
percent law. The source chapter's rounded 70--90 percent description
at printed 641 summarizes part of this modeled phenomenon
\cite[p.~13338]{DeltonEtAl2011Reciprocity}.

The supplement's regressions summarize simulation outcomes, with
adjusted \(R^2=0.456\) for SS1 and 0.573 for SS2, reported coefficient
standard errors and tests. They are not fit to human observations.
No confidence interval for the selected 87/96 percent pair is printed;
none is invented here. The supplement identifies low-benefit/short-
interaction regions where one-shot cooperation does not evolve, and
notes that cooperation in repetition may itself be unfavorable there.
Accidental-defection runs are sensitivity checks within this model,
not independent empirical replications. The primary's Java source was
offered on request, not reproduced in the held paper; no code request
or new simulation run was made
\cite[p.~13336 n.; SI~pp.~2--3, 5--9]{DeltonEtAl2011Reciprocity}.

\subsection{The analytic dependency, with its assumptions}

\begin{hbdefinition}{def:p11-rat106-analytic}{The supplement's fixed-partner payoff problem}
In the no-error analytic model let \(P\in(0,1)\) be the one-shot
base rate, \(0\leq w<1\) the continuation probability of a repeated
dyad, and \(b>c>0\). The common classification probabilities are
\(\Theta_{\mathrm{OS}}\), the probability of a one-shot belief given
an actually one-shot dyad, and \(\Theta_R\), the probability of that
belief given an actually repeated dyad. Each lies in \([0,1]\).
Given the actual dyad type, the two agents' classifications and
randomized action choices are independent. Player \(i\)'s probability
of choosing the cooperative strategy is \(\alpha_i\) given a
one-shot belief, and \(\beta_i\) given a repeated belief.

Cooperation selects TIT-for-TAT; defection selects always-defect.
In an error-free repeated dyad the resulting total-payoff pairs are:
mutual cooperation yields \((b-c)/(1-w)\) to each; cooperative
strategy against always-defect yields \(-c\) to the first and
\(b\) to the second; mutual defection yields zero. The latter two
strategies defect on all rounds after the first. \(V_1\) is player
one's expected payoff before observing its belief classification.
The partner's parameters and the classification probabilities are
held fixed in a best-response calculation.
\end{hbdefinition}

\begin{hbtheorem}{thm:p11-rat106-best-response}{proved=annex:p11.rat106-best-response}{A fully specified one-shot-belief cooperation bridge}
Put, for this theorem's algebra only,
\[
 s_i=\Theta_{\mathrm{OS}}\alpha_i+(1-\Theta_{\mathrm{OS}})\beta_i,
 \qquad r_i=\Theta_R\alpha_i+(1-\Theta_R)\beta_i.
\]
Then the supplement's payoff, equivalently expressed in these conditional
action probabilities, is
\[
 V_1=P(bs_2-cs_1)+(1-P)
       \left(br_2-cr_1+\frac{w(b-c)}{1-w}r_1r_2\right).
\]
It is affine in \(\alpha_1\), with coefficient
\[
 \begin{aligned}
 \frac{w(b-c)}{1-w}(1-P)\Theta_R
       &\bigl(\Theta_R\alpha_2+(1-\Theta_R)\beta_2\bigr)\\
       &{}-c\bigl(P\Theta_{\mathrm{OS}}+(1-P)\Theta_R\bigr).
 \end{aligned}
\]
If this coefficient is positive, the unique best response on
\([0,1]\) is \(\alpha_1=1\); if negative it is
\(\alpha_1=0\); if zero every value ties. With positive
\(w,(1-P),\Theta_R,r_2\), sufficiently large \(b-c\)
makes the coefficient positive. Without those positive factors that
conclusion is not assured. The expression at \(w=0\) remains
well-defined; the source's equivalent \(1/((1/w)-1)\) notation
requires \(w>0\).
\end{hbtheorem}

This supplies the actual payoff bridge in the supplement's S1--S2,
including the equality and degenerate-factor boundaries. It is a
fixed-partner best response, not an evolutionary-convergence theorem.
For instance, with \(\Theta_R=0\) and
\(\Theta_{\mathrm{OS}}>0\), the coefficient is negative:
certainly identifying a one-shot dyad restores the material-payoff
incentive to defect. With \(w=0\) there is no future cooperative
surplus. Conversely a posterior classification as most likely one-shot
can coexist with a positive expected payoff from cooperation when
the less likely repeated opportunity is sufficiently valuable.
Acting optimally under uncertainty need not mean acting as if the
most likely state were certain
\cite[SI~pp.~3--4]{DeltonEtAl2011Reciprocity}.

The justified simulation conclusion is existence and robustness of
one-shot cooperation within these parameterized architectures under
individual-level selection. Group selection is not needed to generate
that modeled result; it is not thereby refuted as a possible contributor
to real cooperation. The main paper itself does not claim direct
reciprocity is the only force shaping human cooperation. Transfer from
this model to observed generosity requires evidence about actual cues,
beliefs, payoffs, recognition, strategies and transmission. A laboratory
instruction is a cue whose reception must be studied, not automatically
the model's perfect knowledge. Likewise, a fitness-favored tendency does
not establish a moral obligation to cooperate or an epistemic entitlement
to believe an interaction repeats.

\section{Comparisons, objections and consultation}
\label{sec:p11-rat106-comparison}

Ecological rationality in Chapter~\ref{sec:p17-rationality-8-5}
evaluates procedures against environmental structure and stated goals,
including tools acquired individually, culturally or through evolution.
The present chapter shares its demand to model inputs and task fit,
but specifically proposes selected ancestral function as the benchmark
for an evolved adaptation. A newly learned heuristic can work well in
a present environment without having that ancestral function; an
ancestrally useful mechanism can fail in today's environment. Present
success and evolutionary origin are not interchangeable explanations.

Evolutionary game theory in
Chapter~\ref{sec:p17-rat93-evolutionary-games} provides explicit
transmission and population models. The adaptationist chapter uses it
to hypothesize adaptive tasks and cognitive design. Neither an ESS,
a rest point nor strategy prevalence is itself evidence of accurate
belief or moral merit. The P17 treatment also distinguishes biological
reproduction from learning and cultural revision, and shows why
stability, convergence and welfare have different conditions. A
generic verbal appeal to selection cannot replace that model.

The epistemic account in
Section~\ref{sec:p11-ri-theoretical-practical} asks how beliefs relate
to truth-directed evidence, while the present functional account can
assess a biased classification by its asymmetric consequences. SS1
preserves Bayesian belief formation but changes the action mapping;
SS2 changes the classification rule. Neither case logically shows
that truth-directed standards are false. The disagreement concerns
what is evaluated and which good is being used as the standard.
Chapter~\ref{sec:p11-rationality-2-2} further separates belief,
choice, mental rationality and external correctness; adopting a common
word does not supply their unification.

Four objections constrain the proposal. First, a plausible ancestral
story may underdetermine the selected function: the same output can
arise through learning, domain-general inference or a different adaptation.
Predictions distinguishing these alternatives are necessary. Second,
task improvements do not automatically establish evolution: a format
effect can support a representation claim while leaving genealogy open.
Third, fitness and epistemic accuracy can diverge, especially where
error costs are unequal; deciding which constitutes rationality needs
the explicit evaluative bridge. Fourth, contemporary tasks can genuinely
require acquired formal methods. Explaining a mismatch with ancestral
inputs does not solve that present problem or excuse a dangerous error.
These are methodological and normative limits of the reconstructed
argument, not claims that every cited empirical study has been refuted.

\begin{hbmethod}{def:p11-rat106-consultation}{Choosing and testing the intended benchmark}
State whether the question concerns evolved design, present task
success, truth-directed belief, preference coherence or moral obligation.
Identify the candidate function and distinguish its hypothesis from
the mechanism's existence. Specify available cues, information format,
costs of the different errors, temporal horizon, alternatives and
transmission assumptions. Derive what the chosen benchmark actually
predicts before applying an error label.

Test a contrast that could separate explanations: natural versus
normalized counts; need below versus above the sure outcome; absent
versus present activation cues; logical violation versus exchange
cheating; or known-one-shot versus uncertain repetition. Preserve
the formal assumptions and the participants' understood task.

For empirical conclusions construct the record in
Method~\ref{def:p19-evidence-record}; a survey percentage without
its population and design is not an independent intervention estimate.
For simulation conclusions identify the architecture, parameters,
stochastic replication and what was measured, then state the missing
real-world mapping. For normative conclusions give the evaluative
bridge independently of frequency or reproductive advantage.
\end{hbmethod}

The orchard count succeeds as a probability calculation but cannot
establish the evolutionary origin of frequency processing. The disease
threshold succeeds as a stated criterion but changes if the relevant
goal is instead another utility. Exchange-content search can succeed at
its agreement-enforcement goal while answering a different logical
question. The reciprocity best response succeeds only with the declared
continuation, cue and partner conditions. These successful and failing
applications express the source's meta-normative contribution without
converting adaptation into an unrestricted endorsement of human thought.

\index{adaptationism!meta-normative rationality}
\index{rationality!adaptive function}
\index{rationality!ecological}
\index{reasoning!social exchange}
\index{cooperation!uncertainty about repetition}
\index{frequency formats}
